\scanpage{130}
% archive sheet number: 129

Supposons maintenant que \deleted{$H,H'$} l'on
se donne une alg. Azumaya $A$ et des
plongements

\margin{plus \unsure{gén.|génl}, dans une $A$ admettant \unsure{formules|formule|fournissant} H.C. d'où $\Delta_A,\mathrm{Tr}_A$ .....}

\diagram{\begin{tikzcd}[ampersand replacement=\&,row sep=small]
\& A \& \\
H \arrow[ur,"i"] \& \& H' \arrow[ul,"j"']
\end{tikzcd}}

On peut alors définir hom de $k$-modules

\diagram{\begin{tikzcd}[ampersand replacement=\&]
P \arrow[r,"\deleted{\unsure{t}}f"] \& A
\end{tikzcd}}

par la condition

\[
 f(\deleted{\unsure{x}}u\otimes v)=i(u).j(v)
\]

\deleted{ce qui induit}
et on récupère
\deleted{et la donnée de}
$i,j$ par

\[
 i=f\circ i_0,\quad j=f\circ j_0
\]

La donnée de $i,j$ équivaut à la
donnée de $f$, soumis aux
conditions que $f\mid H$ et $f\mid H'$
soient des hom d'algèbres. On a alors

\diagram{\begin{tikzcd}[ampersand replacement=\&,row sep=small,column sep=small]
\& A \& \\
\& P \arrow[u,"\unsure{f|t}"'] \arrow[d] \& \\
\& B \arrow[dl] \arrow[dr] \& \\
H \arrow[uur, crossing over, "\deleted{\nearrow}" description,near end] \& \& H' \\
\& k \arrow[ul,"\varepsilon_H"] \arrow[ur,dash,"\varepsilon_{H'}"'] \&
\end{tikzcd}}

\margin{attention, on ne suppose pas que $f$ soit un iso de Mod i.e que \deleted{\unsure{$A$}}$H,H'$ engendrent $A$ comme algèbre}

